\subsection{Path B: Literature-Calibrated Microsimulation}\label{sec:pathb}

\subsubsection{Specification}

The first hazard path brings competing risks, literature-calibrated elasticities, and a cleaner isolation of the Danish counterfactual, without reintroducing the ABM's utility-maximizing machinery. It simulates a stratified sample of 75,000 loans, 40,234 of them still active when the QT window opens (Appendix~\ref{app:params}). Each loan carries static covariates---FICO, region, original LTV, coupon, and a stratum identifier keyed on the Path A specification less its LTV axis---and dynamic state: loan age, balance, and a regime-specific rate gap. Equity enters only through the static original-LTV covariate: no updated CLTV and no local house-price feed exists anywhere in the pipeline, so equity dynamics over the window are unmodeled, a disclosed design boundary for this turnover-dominated, deeply out-of-the-money window. The FICO/LTV covariate ablation reported below brackets what the static block contributes. I specify each loan's monthly hazard of prepayment as
\begin{equation}\label{eq:pathB}
h^{\mathrm{prep}}_{i,t} = \max\!\left\{\, \underline{h},\; h_0\!\left(a_{i,t}\right) \exp\!\left( \beta_1\, g_{i,t} + \beta_F\, z^{\mathrm{FICO}}_{i} + \beta_L\, z^{\mathrm{LTV}}_{i} + \beta_B\, B_{s(i),t} \right) \right\},
\end{equation}
with a competing hazard of default
\begin{equation}\label{eq:default}
h^{\mathrm{def}}_{i,t} = h^{\mathrm{def}}_0 \exp\!\big( \gamma_r\, \sigma_{i,t} + \gamma_F\, z^{\mathrm{FICO}}_{i} + \gamma_L\, z^{\mathrm{LTV}}_{i} \big), \qquad \sigma_{i,t} = \max\{0,\ -\tilde g_{i,t}\},
\end{equation}
where $\sigma_{i,t}$ is rate stress, $h^{\mathrm{def}}_0 = 3\times10^{-4}$ monthly, and $\tilde g_{i,t} = g_{i,t}/100$ is \eqref{eq:pathB}'s rate gap in decimal units, so $\sigma$ is the decimal magnitude of an out-of-the-money gap. The default intercept is a prior implying $0.36\%$ per year, against a realized Freddie Mac default rate of $0.31\%$ over the window on this paper's own GMAR extraction; the prior was fixed before that series was assembled, so the calibration did not use it. Every Path B parameter is fixed a priori, and none is estimated on the benchmark. Appendix~\ref{app:params} collects the parameter vector with provenance (Table~\ref{tab:params}), the delinquency transition matrix, the production step order, and the competing-risks taxonomy; Table~\ref{tab:hazard-notation} maps every symbol to its Path A counterpart.

Here $g_{i,t}$ is the rate gap---coupon minus prevailing market rate, in 100 bp units, negative when the borrower is locked in---so the positive $\beta_1$ of \eqref{eq:beta1} suppresses prepayment exactly when the par-payoff penalty binds. $B_{s(i),t} \in [0,1]$ is the cohort-level burnout state of loan $i$'s stratum $s(i)$: the stratum's cumulative prepaid UPB as a share of its original UPB, capped at one, every stratum initialised at zero at the window open. It therefore accumulates only within-window prepayment, against an original-UPB denominator that still includes pre-window-terminated loans---window-cumulative, not life-cumulative---and with $\beta_B = -0.5$ it scales the hazard by $e^{-0.5 B} \in [0.61, 1]$. Burnout is deliberately cohort-level; Appendix~\ref{app:params} gives the construction and why a per-loan fractional-burnout measure is not well-defined in this design. The baseline hazard $h_0$ follows the 100~PSA seasoning ramp (CPR rising 0.2 percentage points per month of loan age to 6\% at month 30), applied at monthly frequency via $\mathrm{SMM} = 1-(1-\mathrm{CPR})^{1/12}$.

Finally, $\underline{h}$ is the baseline turnover floor: the monthly single-month-mortality (SMM) equivalent of 4\% annual CPR, $\underline{h} = 1-(1-0.04)^{1/12} \approx 0.0034$. It floors total turnover, not a strictly involuntary component of it: deep-discount-cohort turnover at this level includes discretionary life-cycle moves and cash-out refinancings alongside strictly involuntary events. It is also held \emph{flat in loan age}, an assumption the very ladder used to calibrate it does not support---pooled 2017--2019 out-of-the-money turnover rises monotonically with age, from 3.68\% below twelve months to 10.99\% at twenty-four to thirty-six (Appendix~\ref{app:floormech}). I flag the assumption; I do not sweep it. The reason is that appendix's second finding: the 10.99\% cell is carried entirely by the 2019 falling-rate leg, so the age profile is substantially refinancing-contaminated, and an age-varying floor built from it would inherit that contamination. An age dimension therefore remains the one calibration axis this paper sweeps in neither level nor time; a clean version of it needs a rate-regime-pure age ladder this design does not have. No component decomposition exists in this design either, and the strictly-involuntary share (death, divorce, and forced relocation) is plausibly well under half. The floor's semantics do institutional work in selecting the U.S.-intercept Danish anchor (Section~\ref{sec:abm-danish}) and in retaining the max floor form (Section~\ref{sec:robustness-floor}), where the operative premise is only that baseline turnover cannot fall below its observed deep-discount level when the payoff rule changes---which the total-turnover reading supports. Those choices are conventions consistent with that premise, not consequences of a strictly-involuntary decomposition.

The floor's level is anchored to the same empirical observation as the ABM's mobility calibration (Appendix~\ref{sec:method-frictions}): turnover on deeply out-of-the-money discount cohorts during 2023--2024. That period lies \emph{inside} the June 2022--November 2025 evaluation window. No Path B parameter touches the benchmark dollar figure, but the specification's dominant level-setter is measured from in-window prepayment behavior, so the verification content is the marginal over the $\beta_1 = 0$ null, not the aggregate level itself. The pre-committed floor sweep of Section~\ref{sec:robustness-floor} shows that even this marginal is conditional on the floor calibration across the empirically plausible 3--5\% floor range (Table~\ref{tab:floorband}). That interval, not any point inside it, is the identified object, and its positivity at every floor and elasticity-band point swept is a consistency check the window's rate configuration cannot let fail (Section~\ref{sec:identification}). The floor's functional form matters too. Replace the hard maximum of \eqref{eq:pathB} with a competing-risks combination, under which the elasticity is never censored, and the production marginal moves to $+11.25$ points and turns nearly floor-invariant ($+11.29$ to $+11.21$ across the same 3--5\% range). The max-form marginal's floor-dependence is censoring-driven; it is not evidence about the elasticity. The floor binds in 36\% of the 1,683,124 evaluated loan-months in the production U.S. run, so a substantial share of the aggregate level rides on it. That is consistent with the no-lock-in null of Section~\ref{sec:hazard-interp}, where the rate-inelastic model recovers 97.8\% of the benchmark under the standalone scorer (88.7\% on the shared accounting basis). Those are production-floor-form levels; under the additive form the same null recovers 44.7\% at the off-window floor against the production form's 94.8\%, both on the standalone scorer. Proportionality in \eqref{eq:pathB} accordingly holds above the floor, not through it.

The rate-gap coefficient is calibrated from \citet{liebersohn2024}, who estimate that a 100 bp increase in the prevailing rate above the locked-in coupon reduces quarterly mobility probability by 5.5\% to 7.7\% depending on specification. I adopt the 6.5\% midpoint of that band as the ex-ante central scenario, with the band edges as sensitivity bounds, converted to a monthly hazard via the survival-function transform
\begin{equation}\label{eq:beta1}
\beta_1 = -\ln\!\left[\frac{1-\big(1-(1-\delta)\,P_q\big)^{1/3}}{1-\big(1-P_q\big)^{1/3}}\right],
\end{equation}
where $\delta$ is the proportional quarterly mobility decline per 100 bp of rate gap, and $P_q = 0.06$ a reference baseline turnover probability occupying the transform's quarterly slot. That reference is a conversion auxiliary, not a measured quarterly mobility rate: it sits well above \citepos{liebersohn2024} own 1.5\% zero-gap quarterly moving level (Section~\ref{sec:robustness-crossdesign}), and I retain it because quadrupling the assumed moving level over $P_q \in (0, 0.12]$ moves the imported coefficient by under four percent (Appendix~\ref{app:params}, which also documents the superseded continuous-hazard variant, a distinct object the simulation does not use). Equation~\eqref{eq:beta1} is the transform implemented in the codebase; at the central $\delta = 0.065$ it evaluates to $\beta_1 = 0.0686$, rounded to $0.069$ in Table~\ref{tab:params}. It yields the dimensionless log-hazard-ratio implied by a quarterly probability decline; it replaces a divide-by-three approximation, not a hazard itself. Then \eqref{eq:pathB} applies $\beta_1$ linearly per 100 bp of gap. At QT-era gaps of roughly $-380$ bp this compounds the local quarterly estimate to a hazard ratio of about $0.77$ ($\approx e^{3.8 \times -0.069}$), a log-linear extrapolation beyond the range in which the source elasticity was estimated. I flag it as a functional-form assumption and test it with a pre-committed concave-transform rerun of the paired legs: the extrapolation proves immaterial to the level and material to the marginal on the relative reading that ranking rests on. The transform dimension travels with the marginal's operative uncertainty (Table~\ref{tab:uncertainty}; the test and its materiality reading are in Appendix~\ref{app:params}).

\subsubsection{Transport assumptions}

Two transport assumptions travel with the import. The first is an aggregation transport that assumes response homogeneity: the source estimand is a ZIP-code-level moving hazard, applied here as a loan-level log-hazard coefficient. The second concerns survival. The estimand is estimated on a population not conditioned on mortgage survival, while \eqref{eq:pathB} applies it to a pool from which 34,734 of the 75,000 sampled loans had already prepaid before the window opened---a survival-selected pool depleted of the rate-responsive. This transport's direction is therefore signed toward attenuation, and priced: under the gamma-frailty selection identity the attenuation is $\eta = S_{\mathrm{pre}}^{\xi}$ at the measured pre-window survival share $S_{\mathrm{pre}}$, and the off-window marginal then reads $+5.1$/$+4.6$/$+3.7$ points at $\xi = 0.25/0.5/1$. The gamma-frailty parameter $\xi$ is not estimable in this design, and I display it as the free parameter. Path A's burnout coefficient, the one object that speaks to it, is not distinguishable from zero under either bootstrap scheme.

A second mapping assumption governs what the elasticity applies to: \citepos{liebersohn2024} estimand is a moving probability, while \eqref{eq:pathB} applies the implied $\beta_1$ to the entire voluntary hazard, moving and refinancing alike. A pre-committed mixture bracket prices the convention at both floors: $1-\mu$ of the voluntary hazard at its no-response level, $\mu$ at the full response, turnover floor unscaled, parity at $\mu = 1$. The marginal is $+\$37.7$ billion ($+4.9$ points) at $\mu = 0.5$ and $+\$19.2$ billion ($+2.5$ points) at $\mu = 0.25$, against the production $+\$70.3$ billion at $\mu = 1$, close to proportional in $\mu$. At the headlined 4.991\% off-window floor it is $+\$25.7$ billion ($+3.36$ points) at $\mu = 0.5$ and $+\$13.9$ billion ($+1.82$ points) at $\mu = 0.25$, against $+\$42.6$ billion at $\mu = 1$---mildly super-proportional, since the higher floor censors the elasticity in more loan-months and censored contributions do not scale with $\mu$. For this book's rate configuration the convention is therefore signed: applying the imported elasticity to less of the hazard can only shrink the marginal, roughly by its share. No termination-reason field exists in the ingested data \citep[cf.][on the refinancing-versus-moving mix]{amromin2020}, so $\mu$ is a bracketing parameter rather than an estimate, and the production $\mu = 1$ remains the reported form. \citet{graybill2026} corroborate the imported band on exactly this margin: their transaction-record sale-hazard estimate measures the moving component that $\mu$ varies.

The floor enters \eqref{eq:pathB} as a single constant $\underline{h}$ broadcast to every loan-month, with no rate-gap, calendar-time, or covariate argument. Baseline turnover is therefore treated as acyclical, and every rate-responsive movement in prepayment is assigned to the elastic channel. The marginal-over-null design cancels the floor's level---the central and $\beta_1 = 0$ legs carry the identical constant---but it does not cancel a floor co-varying with the rate cycle. To the extent that lock-in also suppresses involuntary moves, that co-movement is absorbed into the estimated elastic marginal, so the in-window marginal is best read as total rate-attributable prepayment suppression, an upper bound on the strictly-voluntary elasticity's contribution. At the off-window calibration the ordering reverses. The demotion's window component (roughly \$20.9 billion; Section~\ref{sec:robustness-floor}) is baseline turnover that the off-window floor classifies as rate-suppressed, and the marginal no longer counts it. So the off-window $+\$42.6$ billion bounds total rate-attributable suppression from below; the matched-depth $+\$63.5$ billion restores that component. The two calibrations bracket two estimands. They are not conservative and liberal readings of one estimand. A within-run cyclical variant sets $\underline{h}_t = 4\% \times (1 + \kappa z_t)$, with $z_t$ the window-standardized 30-year mortgage rate; $\kappa = 0$ nests the production specification and in the run reproduces it exactly. That variant moves the marginal between $+7.1$ and $+9.4$ points across $\kappa \in [-0.5, +0.5]$, every cell inside the level sweep's bracket. Appendix~\ref{sec:robustness-seasonalfloor} decomposes that grid with a pre-committed permutation run: most of its span survives losing the floor path's time order, so the grid measures floor dispersion before it measures cyclicality. The surviving co-movement component is real, small, and bounded there. The window-mean invariant fails, via the floor's non-negativity clip, in the stress cell that sets the bracket's low end.

\subsubsection{Aggregation and the bounds that travel with it}

Aggregation scales simulated prepaid UPB by the ratio of the Fed's MBS book to the 75,000-loan sample balance. Nothing is reweighted toward Ginnie Mae collateral (20.4\% of book face) or vintages outside 2017--2021; the Freddie-estimated hazards extrapolate to those segments as-is. This is not behaviorally one-to-one: Ginnie Mae pools are FHA/VA collateral, with a distinct borrower composition and a servicer buyout channel absent from the conventional book, whose exercise turns on servicer balance-sheet capacity as well as borrower behaviour \citep{aiello2022}. The framework bounds the composition error of treating the Ginnie share as conventional; it does not model that error. Along the coupon dimension, Section~\ref{sec:robustness-hybrid}'s full-book reweighting (sample WAC 3.9\% mapped to the book's 2.49\%) moves Path B from 107.0\% to 109.1\%, about two percentage points---a mixed-basis figure, since 3.9\% is an unweighted mean of borrower note rates and 2.49\% a face-weighted mean of security pass-through coupons. Converted to one convention, the move is $+0.2$ points. Along the agency dimension, Ginnie Mae fixed-rate speeds ran 1.2--2.8 percentage points above Freddie Mac's across the window (5.2\% versus 4.0\%, February 2023; 10.3\% versus 7.5\%, May 2025; Ginnie Mae Global Markets Analysis Report, Recursion data). The involuntary buyout channel is a large minority of that difference, not most of it: the forty-two-month mean Ginnie--Freddie CPR gap of $+2.14$ points splits into a voluntary CRR (conditional repayment rate, prepayment net of the default-driven channel) $+1.35$ (63\%) and a buyout CDR $+0.85$ (40\%), the shares summing past one hundred because the CPR $\approx$ CRR $+$ CDR extraction carries a small residual, and the most buyout-favourable snapshot (CDR 2.1\% versus 0.4\%, May 2025) is the buyout-dominated month, not the representative one. Treating the 20.4\% Ginnie face share as conventional thus mis-states book-wide CPR by at most $0.204 \times (1.2\text{--}2.8) \approx 0.24$--$0.57$ percentage points. The null-to-central sensitivity is roughly \$83 billion per CPR point---window-average book face times one point of annual prepayment times three and a half years---giving a bound of \$20--\$47 billion (2.6--6.2\% of benchmark), signed toward overstating trapped liquidity. The bound is one-sided because the observed differential is: published speeds run Ginnie above Freddie throughout, so the Ginnie-slower branch does not arise in this data.

The vintage dimension carries no bound of either kind. The 2022 vintage alone is 23.1\% of book face, and pre-2017 vintages another 10.6\% (Appendix~\ref{sec:robustness-wal}), both outside the 2017--2021 Freddie sample. The 2022 vintage's lock-in state is not qualitatively different from the sampled vintages', though: of its face, 72.2\% carries a pass-through coupon below 3.0\% (SOMA CUSIP joint coupon $\times$ vintage cells), predominantly late-2021/H1-2022 production four to five points below the window's market rates, and only 13.2\% sits at 4.0\% or above. What the extrapolation lacks is a different \emph{seasoning} state, and neither bound above constrains that. An observed-speed bound in the same template does. Per-quarter cohort cells from the Fannie Mae replication ingest carry observed QT-window prepayment for the out-of-sample vintages: over the window the 2022 vintage prepaid at 4.48\% against the sampled universe's 4.30\% ($+0.17$ unrounded), and the pre-2017 vintages at 5.26\% ($+0.96$; the 2015--2016 cohorts are 99.99\% of that tail's exposure and reproduce its speed unchanged). The share-weighted book-CPR error is $+0.14$ points, or \$11.7 billion at the same \$83-billion-per-CPR-point sensitivity---1.5\% of benchmark, below the static Ginnie bound's own lower edge, yet more than a quarter of the headline lock-in marginal (27.6\% of the off-window \$42.6 billion; a sixth, 16.7\%, of the in-sample \$70.3 billion), signed the same conservative direction: out-of-sample vintages prepaid \emph{faster}, so trapped liquidity is if anything overstated. Applied as a marginal overlay, the template confines the identified marginal to the in-window-vintage share: $+3.7$ points. The overlay scores the 33.7\% out-of-window vintage share by observed Fannie speeds in \emph{both} legs, so that share contributes zero marginal, and the result is pure share scaling ($0.66\times$ against a nominal retained share of $0.663$, placebo within $0.0002$), the vintage counterpart of the Ginnie overlay's $+4.4$. The two overlays cannot be composed without the agency-by-vintage cross-tabulation, since 2.4\% of book face is the 2022--Ginnie intersection. Three caveats travel with the bound: Fannie observed speeds proxy the book's; the pre-2017 leg is a seasoned acquisition tail, not a random draw of old originations; and the 2022 leg is a full-year speed applied to an H1-tilted face.

I extend the static Ginnie bound to a time-varying overlay built from the published speed series; Appendix~\ref{app:params} gives its construction, and its book-wide image enters the assembly of Table~\ref{tab:assembly}. The pipeline also carries two Danish counterfactual mechanisms, only one live in production; the same appendix specifies both, and the anchor that selects between them (results in Section~\ref{sec:abm-danish}).

\subsubsection{Simulated recovery and the calibration band}

Run on the 40,234 window-start survivors (attrition accounting in Appendix~\ref{app:params}), this specification recovers \$818.5 billion, or 107.0\%, of the empirical benchmark. Raw CPR correlation against the empirical path is $r = +0.190$ at zero lag, and peak cross-correlation $r = +0.404$ at lag $-3$. Both are descriptors of the published series, not credentials: four clip-induced zero months dominate the monthly benchmark's differenced signal, so it cannot adjudicate timing (Appendix~\ref{sec:robustness-seasonalfloor}). A CUSIP-level factor-paydown series would remove those artifacts at the source and is the institutionally superior path for future work; the benchmark retains current-face differencing because factor data were not staged in this pipeline. The cross-correlation at lag $k$ pairs the simulated series at month $t$ with the empirical series at month $t+k$, both demeaned over the overlap and with boundary months dropped, and a direction check on synthetic data confirms that a peak at $k=-3$ means the simulated path \emph{trails} the empirical path by three months: the empirical SOMA-derived CPR path moves approximately three months ahead. That peak was selected across seven examined lags, so its nominal significance is overstated, and no inference rests on this single coefficient. The definition was inverted in print, but the direction claims were always computed under the implemented convention; the opposite reading survives in a sha-pinned replication artifact's field, superseded and retired, and the replication package's ledger catalogues it together with the superseded figure and its provenance. The alignment therefore does \emph{not} show the simulated path leading SOMA cash flow. A settlement-delay mechanism---in which loan-level prepayment decisions show up in SOMA cash two to three months later---predicts the \emph{opposite} direction of offset, so I attach no settlement interpretation to it; Appendix~\ref{app:patha-detail} carries the direct input-timing diagnosis behind that reading. The level and band follow the rate-gap coefficient's unit convention of Appendix~\ref{sec:robustness-pipeline}.

The rate-gap elasticity imported from \citet{liebersohn2024} never touches the \$764.7 billion benchmark during calibration. Path A's rate-gap and burnout coefficients, by contrast, are estimated directly from the Freddie Mac sample. The two are mutually reinforcing, not redundant: Path A fits the same suppression channel on the same loans ($+0.29$ per 100 bp; Section~\ref{sec:patha}), and the Fannie Mae replication reproduces its sign on a second agency's loan book at roughly $3.6\times$ the per-basis-point magnitude---directionally consistent with the imported band, but no more. A pre-committed test of $H_0\!: \beta_g \le 0$ respecting the diagnosed bootstrap bias does not reject at any conventional level (BCa, 296-refit jackknife, $p = 0.093$ stratum, $0.412$ temporal; 200-draw month-permutation, $p = 0.241$). The in-sample estimates are therefore not statistically distinguishable from zero under bias-respecting constructions, and the elasticity's evidential basis is the external literature alone. \citet{graybill2026} independently corroborate the band's location: a transaction-record sale-hazard response of roughly 7--11\% per percentage point of rate gap, in a housing-market-equilibrium setting touching no credit file. That corroboration comes from at or above the adopted band---only the 7\% low end falls inside 5.5\%--7.7\%, and the midpoint sits above the high edge. A prepayment-hazard coefficient is also not unit-comparable to \citepos{liebersohn2024} quarterly-mobility decline, and Path A measures prepayment, not pure mobility. What remains is directional consistency, not triangulation of sign or direction, not magnitude. A second anchor corroborates from outside the band. \citepos{fonseca2024} roughly 9\% moving-rate reduction per 100 bp, propagated through the same transform ($\beta_1 \approx 0.096$) and band machinery, yields $+\$87.9$ billion ($+11.5$ points) at the production floor. That figure sits above the Liebersohn--Rothstein high edge ($+\$79.5$ billion, $+10.4$ points), yet inside the pre-committed floor-and-band box, so the central calibration is conservative among the available literature estimates and against the \citet{graybill2026} anchor as well. That estimand is a moving-rate reduction, not a quarterly-mobility decline, so I read it as directional corroboration, not a competing point estimate.

Path B is accordingly a structural verification exercise, not a second fit to my benchmark. The aggregate level is floor-dominated, so landing within 8 percentage points of the benchmark is not by itself evidence for the imported elasticity; the defensible verification claim is the bounded interval of the calibration box (Table~\ref{tab:floorband}), positive at every floor, band, and floor-form point tested. Under this window's rate configuration that positivity verifies the specification. It does not identify the mechanism (Section~\ref{sec:identification}); I do not claim the level validates the elasticity's magnitude. Nor is timing part of that evidence: a zero-elasticity rerun reproduces the same three-month lead (Section~\ref{sec:hazard-interp}), so the lead cannot discriminate the elasticity from the rate-inelastic model.

Propagating the full Liebersohn--Rothstein 5.5\%--7.7\% sensitivity band through the simulation yields \$809.9 billion to \$827.7 billion, or 105.9\%--108.2\% of benchmark under the standalone scorer and 96.8\%--99.1\% on the shared accounting basis of Section~\ref{sec:robustness-hybrid}. Peak $r$ runs between $+0.371$ and $+0.423$; given the moving-block intervals above and the seven-lag search, these differences carry no evidential weight. The central 107.0\% estimate should be read as a point within this interval; Figure~\ref{fig:band} plots the band. That band is a calibration interval, not sampling error. The sampling dimension is quantified separately in Appendix~\ref{app:params}, which reports the committed within-stratum loan bootstrap, and the stratum-cluster scheme that widens it, with their caveats. Even at the cluster scheme's wider interval the sampling layer is less than half the width of the floor-read interval of Section~\ref{sec:robustness-floor}, so the operative uncertainty statements for Path B's level and marginal remain the calibration box and floor-form results there. The shared-basis equivalents follow by the common 9.1-point netting, which cancels from the marginal entirely.

\begin{figure}[H]
\centering
\includegraphics[width=0.82\textwidth]{fig6_elasticity_band}
\caption{Path B trapped liquidity across the Liebersohn--Rothstein elasticity band. The quarterly mobility-decline parameter $\delta$ is swept across its published 5.5\%--7.7\% range (central 6.5\% marked), mapping monotonically to recovered trapped liquidity of \$809.9--\$827.7 billion on the standalone-scorer basis (shared-basis equivalents in Table~\ref{tab:bases}) against the \$764.7 billion benchmark line: stronger assumed mobility suppression yields more trapped liquidity, with the peak cross-correlation lag stable at $-3$ across the entire band. Table~\ref{tab:floorband} extends this band across the baseline-turnover-floor grid.}
\label{fig:band}
\end{figure}

\begin{table}[!t]
\centering
\caption{The calibration box: lock-in marginal over the $\beta_1 = 0$ null, in percentage points of the \$764.7 billion benchmark, across the baseline turnover floor and the Liebersohn--Rothstein band $\delta$. The 6.0\% row's level is reached only by the refi-contaminated pooled and 2019 anchors that Section~\ref{sec:robustness-floor} disqualifies; it is retained as a grid point, not a defensible anchor, and the box's lower edge should be read from the defensible off-window rows of Table~\ref{tab:oosfloor}. The bind column is the share of evaluated U.S. loan-months at which the floor lifts the hazard (central elasticity). Every marginal is basis-invariant (the underlying levels are standalone-scorer figures; both bases in Table~\ref{tab:bases}). Pre-committed runs; artifacts catalogued in the replication package's ledger.}
\label{tab:floorband}
\begin{tabular}{lcccc}
\toprule
Floor (annual CPR) & $\delta = 5.5\%$ & $\delta = 6.5\%$ & $\delta = 7.7\%$ & Floor binds \\
\midrule
2.0\% & $+9.7$ & $+11.3$ & $+13.2$ & 3.9\% \\
3.0\% & $+9.3$ & $+10.8$ & $+12.6$ & 12.1\% \\
3.5\% & $+8.9$ & $+10.3$ & $+11.8$ & 21.6\% \\
4.0\% (production) & $+8.1$ & $+9.2$ & $+10.4$ & 36.3\% \\
4.5\% & $+6.7$ & $+7.5$ & $+8.4$ & 53.5\% \\
5.0\% & $+5.0$ & $+5.5$ & $+6.0$ & 69.1\% \\
6.0\% & $+2.1$ & $+2.3$ & $+2.4$ & 88.7\% \\
\bottomrule
\end{tabular}
\end{table}

\subsection{External Replication on Fannie Mae Loan-Level Data}\label{sec:pathb-fannie}

Path B's construction admits a direct external replication on Fannie Mae's single-family loan performance data (24 acquisition quarters, 2017Q1--2022Q4; 17.6 million loans, 826 million loan-months). I stage it into the identical origination/performance layout, so the panel builder, estimator, and stratified 75{,}000-loan design run unchanged on the same 2017--2021 origination universe. The spec was fixed and committed before the run. Zero-balance code 01 maps to voluntary prepayment as for Freddie; codes 06 and 16---non-credit removals with no Freddie counterpart---are censored, and were they all prepayments the worst-case bracket is 0.38\% of removals. The pre-committed gate required a positive Fannie marginal inside Section~\ref{sec:robustness-floor}'s Freddie floor/band envelope. It replicates: central \$821.0 billion (107.4\% of benchmark), $\beta_1 = 0$ null \$754.6 billion (98.7\%), and a lock-in marginal of $+\$66.4$ billion ($+8.68$ points of benchmark) against Freddie's $+\$70.3$ billion ($+9.20$ points)---inside the envelope. But that envelope, the calibration box, spans 2.11 to 13.17 points: an 11.06-point window around a 9.20-point Freddie estimate, admitting a replication 77\% below it or 43\% above. A gate that wide fails only on a sign error, and I do not treat its passing as evidence. The agreement is a different matter: 17.6 million staged Fannie loans run through the same pipeline land 0.52 points from the Freddie point, far more than the gate demanded, and that closeness is the part of this exercise that is reassuring. Every behavioral parameter is fixed a priori and identical across runs, and the two samples are compositionally indistinguishable (population-stability indices at or below 0.018, Appendix~\ref{app:composition}), so the 0.52-point agreement certifies pipeline determinism and cross-agency compositional stability, not external validity---which lives in the sign reproduction and staging inputs, not in the closeness. Path A's fixed-effects Fannie-cohort refit reproduces the rate-gap coefficient's sign, but at about $3.6\times$ the Freddie fit's per-basis-point magnitude and with Section~\ref{sec:patha}'s no-forward-precision holdout character, so it retains aggregate-corroboration status only. I record two ex ante scope notes. Whether Freddie's code 01 absorbs exits Fannie books under 06/16 is not settled by public documentation, so the claim is sign and envelope-conditional marginal consistency, not a hazard-level match. And Fannie's terms prohibit redistribution, so the package ships ingestion and replication code only---no Fannie data and, conservatively, no derived cells. The committed artifact carries headline statistics alone.
